\documentclass[11pt,reqno]{amsart}

%=====================================================================
%  Team A1 -- Role 1 (Optimal Contract and General Carbon Theory)
%  A1 WEEK 3.  The optimal contract makes the flexible load decarbonize,
%  distribution-free.  Builds on role1_week2_theory.tex (Week 2).
%   1. The carbon-optimal admissible contract is surplus-only.
%   2. Decarbonization theorem: adopted => K_R rises, carbon strictly falls,
%      for EVERY distribution (the sign is distribution-free).
%   3. The triple-win under the contract (profit, reliability, carbon).
%   4. Resolution of Week-2's irreducibility: the contract removes the floor
%      and the distributional ambiguity.
%=====================================================================

\usepackage[T1]{fontenc}
\usepackage{lmodern}
\usepackage{amsmath,amssymb,amsthm,mathtools}
\usepackage[margin=1in]{geometry}
\usepackage{microtype}
\usepackage{enumitem}
\usepackage{booktabs}
\usepackage{xcolor}
\usepackage[colorlinks=true,linkcolor=blue!55!black,citecolor=blue!55!black]{hyperref}

\theoremstyle{plain}
\newtheorem{theorem}{Theorem}
\newtheorem{proposition}{Proposition}
\newtheorem{lemma}{Lemma}
\newtheorem{corollary}{Corollary}
\theoremstyle{definition}
\newtheorem{condition}{Condition}
\newtheorem{assumption}{Assumption}
\theoremstyle{remark}
\newtheorem{remark}{Remark}

\newcommand{\KR}{K_R}
\newcommand{\KF}{K_F}
\newcommand{\KM}{K_M}
\newcommand{\KRc}{K_R^{c}}
\newcommand{\KFc}{K_F^{c}}
\newcommand{\KMc}{K_M^{c}}
\newcommand{\KRo}{K_R^{\circ}}
\newcommand{\KFo}{K_F^{\circ}}
\newcommand{\Dbar}{\bar D}
\newcommand{\CM}{C_M}
\newcommand{\CF}{C_F}
\newcommand{\CR}{C_R}
\newcommand{\bp}{\beta+p}
\newcommand{\FKR}{F_{\!\KR}}
\newcommand{\E}{\mathbb{E}}
\newcommand{\PP}{\mathbb{P}}
\newcommand{\eqdef}{\vcentcolon=}
\newcommand{\fossil}{\mathrm{fossil}}

\begin{document}

\title[The Optimal Contract Makes the Flexible Load Decarbonize (Week 3)]
{The Optimal Contract Makes a Flexible Computing Load Decarbonize the Grid:\\
A Shape-Independent Guarantee}
\author{Team A1 --- Role 1 (Optimal Contract and General Carbon Theory)}
\date{Working results, Week~3}

\begin{abstract}
This is the Role~1 Week~3 deliverable, and it is a positive resolution of the Week-2
carbon question. Week~2 showed that under the \emph{naive} dispatch rule---run the
flexible computing load (\textsc{fcl}) on spare fossil whenever it pays---a finite
interior optimum carries an irreducible fossil floor, so the carbon sign is ambiguous
and, in general, not even decidable from the equilibrium cutoffs. Here we change the
contract and the answer changes with it. We study the regulator-imposed
\emph{surplus contract}, which runs the \textsc{fcl} on renewable surplus and
curtails it in fossil hours, and which minimizes the load's fossil footprint
(driving it to zero). Three results follow, and all are positive. \emph{First}, under
the surplus contract the load burns \emph{zero} fossil, so the Week-2 floor is gone by
construction. \emph{Second}---the keystone---whenever the load is adopted the contract
strictly raises renewable investment, $\KRc>\KRo$, and total carbon \emph{strictly
falls}, $\E[\fossil^{c}]<\E[\fossil^{\circ}]$. The force of this is that, unlike the
naive rule, it needs \emph{no condition on the shape of $F$}: under the standing
regularity (the renewable problem is single-peaked) and adoption, carbon falls whatever
the tails or skew of demand. The carbon sign, which Week~2 showed is decided by the
distributional shape under the naive rule, is freed from that shape by the contract.
\emph{Third}, the contract is a clean Pareto improvement on the no-load benchmark:
profit strictly rises and carbon strictly falls, while reliability is held exactly at
the no-load level (the surplus load never touches fossil, so it leaves the blackout
target unchanged). The two standing hypotheses are adoption ($\CM<R\,F(0;\KRo)$, the
load is cheap enough to run on the available surplus) and the Week-2 single-peakedness;
neither restricts the shape of $F$, and we are explicit that the surplus optimum may sit
at the renewable-capacity bound when the surplus channel is strong. Every identity is
checked against quadrature in \texttt{verify\_week3.py}. The message is clean: a flexible
computing load is a decarbonization instrument, and contract design is the lever that
makes it one.
\end{abstract}

\maketitle

%---------------------------------------------------------------------
\section{Setup: the model and the contract space}\label{sec:setup}

We keep the model and notation of the Week-1 and Week-2 deliverables. A regulated
utility chooses renewable capacity $\KR$, dispatchable fossil $\KF$, and
flexible-computing-load (\textsc{fcl}) capacity $\KM$. Residual demand is
$\tilde x=\tilde D-\KR\tilde S\in[-\KR,\Dbar]$ with CDF $F(x;\KR)$, density
$f(x;\KR)>0$ on $(-\KR,\Dbar)$, and $\FKR\eqdef\partial F/\partial\KR>0$ on the
interior of the support (more renewables strictly shift residual demand toward surplus
in every interior state---true for the scale model $\tilde x=\tilde D-\KR\tilde S$,
where $\FKR=\E[\tilde S\,f_{\tilde D}(x+\KR\tilde S)]>0$ wherever $f_{\tilde D}>0$). This
positivity is the only structural property of $F$ we use, and it holds across surplus
states ($x<0$) and fossil states ($x>0$) alike. Per-MWh: \textsc{fcl} revenue $R$,
fossil cost $c$, price $p$, blackout penalty $\beta$; amortized capital costs
$\CR,\CF,\CM$.

\begin{assumption}[Regularity, as in Week 1]\label{ass:reg}
$R>c+\CM$,\quad $R<c+\CF+\CM$,\quad $\bp>c+\CF$,\quad $\bp>R$.
\end{assumption}

\paragraph{Dispatch and fossil accounting.} In state $x$, consumers are served first
(they outrank the \textsc{fcl} for scarce fossil), so consumer fossil is
$\min\{(x)_+,\KF\}$. Renewable surplus $(-x)_+$ is free power. A \emph{dispatch
contract} is a measurable schedule $d:\mathbb R\to[0,\KM]$ with
$d(x)\le(\KF-x)_+$ (the \textsc{fcl} can draw surplus plus fossil headroom); the
fossil it burns is $(d(x)-(-x)_+)_+$ (surplus first, then fossil). Total fossil is
\begin{equation}\label{eq:totfossil}
  \E[\fossil]=\underbrace{\E\!\big[\min\{(\tilde x)_+,\KF\}\big]}_{\text{consumer fossil}}
  \;+\;\underbrace{\E\!\big[(d(\tilde x)-(-\tilde x)_+)_+\big]}_{\text{\textsc{fcl} fossil}} .
\end{equation}
Consumer fossil equals $\int_0^{\KF}(1-F)\,dx$ and does not depend on the contract
$d$; only the second term does.

\paragraph{Two contracts.} The Week-2 analysis used the \emph{firm-optimal} (naive)
contract $d^{*}(x)=\min\{(\KF-x)_+,\KM\}$, which runs the load whenever the marginal
$R-c>0$ is positive---hence on fossil headroom too. We contrast it with the
\emph{surplus contract}
\begin{equation}\label{eq:surpcontract}
  d^{c}(x)\eqdef\min\{(-x)_+,\KM\} ,
\end{equation}
which runs the load only on renewable surplus and curtails it whenever fossil would be
needed. We separate the two decision-makers cleanly: the \emph{regulator} sets the
real-time dispatch contract $d(\cdot)$ (the regulated utility dispatches as directed),
and the \emph{firm} then chooses capacities $(\KR,\KF,\KM)$ to maximize profit under
that contract. The surplus contract is regulator-imposed and is \emph{not} dispatch
incentive-compatible: with $R>c$ the firm would prefer to run the load on fossil
headroom too (the naive $d^{*}$). It is enforced as a curtailment rule, exactly as
demand-response and interruptibility tariffs are. The no-\textsc{fcl} benchmark
$(\KRo,\KFo)$ is the program with $\KM=0$, with fossil
$\E[\fossil^{\circ}]=\int_0^{\KFo}(1-F(x;\KRo))\,dx$ and the distribution-free
cutoff
\begin{equation}\label{eq:phio}
  F(\KFo;\KRo)=\varphi^{\circ}\eqdef\frac{\bp-c-\CF}{\bp-c}.
\end{equation}

%=====================================================================
\section{The fossil-minimizing dispatch is surplus-only}\label{sec:optcontract}
%=====================================================================

We first show that, at any fixed capacities, the surplus contract minimizes the load's
fossil footprint---driving it to zero---and that among zero-footprint contracts it is
the one that monetizes the most surplus, hence the natural carbon-oriented contract.
(We do \emph{not} claim global carbon-optimality over the bilevel game: a different
contract could induce a larger $\KR$ at the price of some \textsc{fcl} fossil. Week~2
showed that trade---the naive rule buys a bigger capacity bump but pays the fossil
floor---loses on net; the surplus contract takes the other side, zero footprint with a
strict capacity gain, and \S\ref{sec:decarb} shows that side wins distribution-wide.)

\begin{proposition}[Carbon-optimal dispatch]\label{prop:surplusopt}
Fix capacities $(\KR,\KF,\KM)$. Among all admissible dispatch schedules $d$, total
fossil \eqref{eq:totfossil} is minimized, state by state, by running the load on
surplus alone, $d^{c}(x)=\min\{(-x)_+,\KM\}$, at which the \textsc{fcl} fossil term is
identically zero. Any schedule that ever dispatches beyond the surplus,
$d(x)>(-x)_+$, burns strictly more fossil in that state.
\end{proposition}

\begin{proof}
Consumer fossil in \eqref{eq:totfossil} is independent of $d$, so minimizing total
fossil is minimizing $\E[(d(\tilde x)-(-\tilde x)_+)_+]$ pointwise. The integrand
$(d(x)-(-x)_+)_+$ is zero whenever $d(x)\le(-x)_+$ and strictly increasing in $d(x)$
beyond that. Setting $d(x)=\min\{(-x)_+,\KM\}\le(-x)_+$ attains zero; any
$d(x)>(-x)_+$ gives a strictly positive value. Hence $d^{c}$ is the state-by-state
fossil minimizer, with \textsc{fcl} fossil $\equiv0$. Among the zero-fossil schedules
($d\le(-x)_+$), $d^{c}=\min\{(-x)_+,\KM\}$ is the one that runs the load to the full
available surplus, so it maximizes \textsc{fcl} revenue $R\,\E[d]$ subject to zero
fossil; it is therefore the carbon-optimal contract that the firm will actually operate
(any $d<d^{c}$ earns less surplus revenue for the same zero footprint, and not
operating at all, $d\equiv0$, forgoes the revenue that makes adoption worthwhile).
\end{proof}

So the surplus contract is the one a regulator who wants to keep the load's
benefits while erasing its fossil footprint would write: monetize the surplus, never
burn fossil to mine. Under it, \eqref{eq:totfossil} collapses to consumer fossil,
\begin{equation}\label{eq:fosc}
  \E[\fossil^{c}]=\int_0^{\KFc}\big(1-F(x;\KRc)\big)\,dx ,
\end{equation}
exactly the no-\textsc{fcl} form---but evaluated at the capacities the contract
induces, which is where the gain comes from.

%=====================================================================
\section{The decarbonization theorem (no shape condition)}\label{sec:decarb}
%=====================================================================

Under the surplus contract the firm's program is
\begin{equation}\label{eq:Pic}
  \Pi^{c}=p\,\E[\tilde D]
  +R\!\!\int_{-\KM}^{0}\!\!\!F\,dx
  -c\!\!\int_{0}^{\KF}\!\!(1-F)\,dx
  -(\bp)\!\!\int_{\KF}^{\Dbar}\!\!(1-F)\,dx
  -\CF\KF-\CR\KR-\CM\KM ,
\end{equation}
the \textsc{fcl} contributing the surplus revenue $R\,\E[\min\{(-\tilde x)_+,\KM\}]
=R\int_{-\KM}^{0}F\,dx$ and \emph{no} fossil cost. Its first-order conditions are
\begin{align}
  \partial_{\KM}\Pi^{c}&=R\,F(-\KM;\KR)-\CM=0
     &&\Longrightarrow\quad F(-\KMc;\KRc)=\tfrac{\CM}{R},\label{eq:focM}\\
  \partial_{\KF}\Pi^{c}&=(\bp-c)\big(1-F(\KF;\KR)\big)-\CF=0
     &&\Longrightarrow\quad F(\KFc;\KRc)=\varphi^{\circ},\label{eq:focF}\\
  \partial_{\KR}\Pi^{c}&=\underbrace{R\!\!\int_{-\KM}^{0}\!\!\FKR\,dx}_{\ge0,\ \text{surplus channel}}
     +\,c\!\!\int_{0}^{\KF}\!\!\FKR\,dx+(\bp)\!\!\int_{\KF}^{\Dbar}\!\!\FKR\,dx-\CR=0.\label{eq:focR}
\end{align}
Two facts jump out. The fossil cutoff \eqref{eq:focF} is \emph{identical} to the
no-\textsc{fcl} one \eqref{eq:phio}: $F(\KFc;\KRc)=\varphi^{\circ}=F(\KFo;\KRo)$. And
the renewable condition \eqref{eq:focR} is the no-\textsc{fcl} renewable condition plus
the strictly positive \emph{surplus channel} term $R\int_{-\KM}^{0}\FKR\,dx$. The
\textsc{fcl} is adopted ($\KMc>0$) when \eqref{eq:focM} has an interior root; a
sufficient condition (using $F(0;\KRc)\ge F(0;\KRo)$, since $\KRc\ge\KRo$ and $F(0;\cdot)$
is increasing) is
\begin{equation}\label{eq:adopt}
  \CM<R\,F(0;\KRo)\qquad\text{(adoption: the load is cheap enough to run on the surplus).}
\end{equation}

\begin{condition}[Single-peaked renewable problem]\label{cond:reg}
The renewable marginal value (the left side of \eqref{eq:focR} plus $\CR$) is strictly
decreasing in $\KR$ along the locus \eqref{eq:focF}. This is the Week-2 reduced
second-order condition (the Schur/diagonal-dominance condition), here for the
surplus program; it holds on the regular region and at the power-law boundary it is
the knife-edge.
\end{condition}

\begin{theorem}[Decarbonization, free of any shape condition]\label{thm:decarb}
Suppose the load is adopted \eqref{eq:adopt} and Condition~\ref{cond:reg} holds. Then
under the surplus contract
\[
  \boxed{\;\KRc>\KRo\quad\text{and}\quad \E[\fossil^{c}]<\E[\fossil^{\circ}]\;}
\]
where $\KRc$ is the surplus optimum (an interior root of \eqref{eq:focR} if one exists,
else the renewable-capacity bound $\bar K_R$; \S\ref{sec:decarb}). The conclusion uses
only the standing positivity $\FKR>0$ on the interior support, Condition~\ref{cond:reg},
and adoption---\emph{no} restriction on the shape, tails, or skew of $F$. This is the
sharp contrast with Week~2, where (Remark~\ref{rem:resolve}) the conditional shape of
$F$ decided the sign; the surplus contract removes that shape-dependence.
\end{theorem}

\begin{proof}
\emph{Step 1 ($\KRc>\KRo$).} Evaluate $\partial_{\KR}\Pi^{c}$ at the no-\textsc{fcl}
optimum $(\KRo,\KFo)$, taking $\KM$ from \eqref{eq:focM} at $\KR=\KRo$ (which is
$>0$ under \eqref{eq:adopt}). By \eqref{eq:focF} the fossil cutoff at $\KRo$ is
$\KFo$, so the last two integrals in \eqref{eq:focR} are exactly the no-\textsc{fcl}
renewable benefit, which equals $\CR$ at $\KRo$. Hence
\[
  \partial_{\KR}\Pi^{c}\big|_{\KRo}
  = R\!\int_{-\KM}^{0}\!\FKR(x;\KRo)\,dx \;>\;0 ,
\]
strictly, because $\FKR\ge0$ with $\FKR>0$ on a set of positive measure in
$(-\KM,0)$ (more renewables raise the surplus probability). So $\Pi^{c}$ is still
strictly increasing in $\KR$ at $\KRo$; by Condition~\ref{cond:reg} the renewable
marginal value is strictly decreasing, so if it reaches zero at a finite $\KR$ that
zero $\KRc$ lies strictly to the right of $\KRo$. If instead the marginal value stays
positive up to the renewable-capacity bound $\bar K_R$ (the surplus channel makes
renewables so attractive that the unconstrained program is a corner---see
Proposition~\ref{prop:saving}), the optimum is $\KRc=\bar K_R$. Either way
$\KRc>\KRo$: the surplus contract strictly raises renewable investment.

\emph{Step 2 ($\KFc\le\KFo$).} Both cutoffs satisfy $F(\cdot;\cdot)=\varphi^{\circ}$:
$F(\KFc;\KRc)=F(\KFo;\KRo)=\varphi^{\circ}$. Since $F$ is increasing in $\KR$
($\FKR\ge0$), $F(\KFo;\KRc)\ge F(\KFo;\KRo)=\varphi^{\circ}=F(\KFc;\KRc)$, and as
$F(\cdot;\KRc)$ is increasing in $x$ this forces $\KFc\le\KFo$.

\emph{Step 3 (carbon strictly falls).} By \eqref{eq:fosc} and
$\E[\fossil^{\circ}]=\int_0^{\KFo}(1-F(x;\KRo))\,dx$,
\[
  \E[\fossil^{\circ}]-\E[\fossil^{c}]
  =\int_0^{\KFo}\!\!(1-F(x;\KRo))\,dx-\int_0^{\KFc}\!\!(1-F(x;\KRc))\,dx .
\]
Write this as
$\int_0^{\KFc}\![F(x;\KRc)-F(x;\KRo)]\,dx+\int_{\KFc}^{\KFo}(1-F(x;\KRo))\,dx$.
The first integral is $\ge0$ because $\KRc>\KRo$ gives $F(\cdot;\KRc)\ge F(\cdot;\KRo)$,
and $>0$ since $\FKR>0$ on a positive-measure subset of $(0,\KFc)$; the second is
$\ge0$ since $\KFc\le\KFo$ and $1-F\ge0$. Hence the difference is strictly positive:
$\E[\fossil^{c}]<\E[\fossil^{\circ}]$. None of the three steps used any feature of $F$
beyond the standing positivity $\FKR>0$ on the interior support (Step~1 uses it on the
surplus side $(-\KM,0)$, Step~3 on the fossil side $(0,\KFc)$) and
Condition~\ref{cond:reg}, so the conclusion holds with \emph{no condition on the shape
of $F$}.\footnote{Verified in \texttt{verify\_week3.py}: across Beta
scale-mixtures (symmetric, left- and right-skewed, bimodal) the surplus-channel term
is strictly positive at $\KRo$, $\KRc>\KRo$, $\KFc\le\KFo$, and
$\E[\fossil^{c}]<\E[\fossil^{\circ}]$ at every adopted optimum---the sign never flips.}
\end{proof}

\begin{remark}[Why this resolves Week 2]\label{rem:resolve}
Week~2 proved that under the naive contract the carbon sign is \emph{irreducibly
distributional}: at fixed equilibrium capacities the admissible conditional laws make
$\Delta$ range over an interval straddling zero, so no distribution-free statement is
possible and a fossil floor $(\KM-\KR)_+$ blocks reduction at interior optima.
Theorem~\ref{thm:decarb} shows the obstruction was an artifact of the \emph{contract},
not the technology. The surplus contract sets the \textsc{fcl} fossil---and with it the
floor---to zero, and the only remaining channel is the renewable build-out, which has a
definite sign. The carbon question goes from ``it depends on the tails of $F$'' to
``it always falls when the load is adopted.''
\end{remark}

%=====================================================================
\section{The triple-win and the exact saving}\label{sec:triplewin}
%=====================================================================

\begin{corollary}[The optimal contract Pareto-dominates no load]\label{cor:triplewin}
Under the hypotheses of Theorem~\ref{thm:decarb}, relative to the no-\textsc{fcl}
benchmark the surplus contract \textup{(i)} strictly raises the utility's profit and
\textup{(iii)} strictly lowers carbon, while \textup{(ii)} holding the blackout
probability \emph{exactly} at the no-load level. It is thus a strict Pareto
improvement: more profit and less carbon at no reliability cost. The contract is
regulator-imposed, not firm-optimal: left to itself the firm would also run the load on
fossil headroom (the naive rule) and earn still more. But it strictly prefers operating
under the surplus contract to having no load, so participation is individually rational
and no subsidy is needed---the regulator gives up only the firm's extra
\emph{fossil-powered} computing profit, which is exactly the carbon-causing margin.
\end{corollary}

\begin{proof}
(iii) is Theorem~\ref{thm:decarb}. (i): the surplus contract adds the option of
earning $R$ on otherwise-curtailed surplus at hardware cost $\CM$; by \eqref{eq:adopt}
the first \textsc{fcl} unit's marginal surplus revenue $R\,F(0;\KRo)$ exceeds $\CM$, so
$\Pi^{c}>\Pi^{\circ}$ strictly (the firm can always replicate the benchmark by setting
$\KM=0$, and adoption is profitable). (ii): blackout occurs on $\{\tilde x>\KF\}$ with
probability $1-F(\KF;\KR)$. By the fossil cutoff \eqref{eq:focF}, \emph{both} programs
target the same reliability, $F(\KFc;\KRc)=\varphi^{\circ}=F(\KFo;\KRo)$, so the
blackout probability is $1-\varphi^{\circ}$ in both---identical. The surplus load never
uses fossil, so it leaves the fossil-capacity trade-off, and hence the blackout target,
untouched; the larger $\KRc$ is spent on \emph{lowering fossil capacity}
($\KFc\le\KFo$, Step~2) rather than on extra reliability. Reliability is therefore held,
not improved---a clean tie, which is what makes the result a Pareto
improvement.\footnote{Verified in \texttt{verify\_week3.py}: $\Pi^{c}>\Pi^{\circ}$ and
blackout$^{c}=$blackout$^{\circ}$ (to numerical tolerance) at every adopted optimum.}
\end{proof}

\begin{proposition}[Exact carbon saving and its monotonicity]\label{prop:saving}
The carbon saving under the surplus contract is
\begin{equation}\label{eq:exactsave}
  \Delta^{c}\eqdef\E[\fossil^{\circ}]-\E[\fossil^{c}]
  =\int_0^{\KFc}\!\!\big[F(x;\KRc)-F(x;\KRo)\big]\,dx
   +\int_{\KFc}^{\KFo}\!\!\big(1-F(x;\KRo)\big)\,dx \;>\;0 .
\end{equation}
It is increasing in the surplus channel: a larger adopted $\KMc$ (cheaper load, more
surplus) raises $\KRc$ and hence $\Delta^{c}$. In the cheap-renewable limit
$\KRc\to\infty$ the contract drives the utility to full decarbonization,
$\E[\fossil^{c}]\to0$, so $\Delta^{c}\to\E[\fossil^{\circ}]$.
\end{proposition}

\begin{proof}
The identity is Step 3 of Theorem~\ref{thm:decarb}. Monotonicity: since
$\Delta^{c}=\E[\fossil^{\circ}]-\E[\fossil^{c}]$ and $\E[\fossil^{\circ}]$ is fixed, it
suffices that $\E[\fossil^{c}]=\int_0^{\KFc}(1-F(x;\KRc))\,dx$ is \emph{decreasing} in
$\KRc$ (a larger adopted $\KMc$ enlarges the surplus channel and raises $\KRc$ by
Condition~\ref{cond:reg}). Differentiating, with $1-F(\KFc;\KRc)=1-\varphi^{\circ}$ at
the moving upper limit,
\[
  \frac{d\,\E[\fossil^{c}]}{d\KRc}
  =\underbrace{(1-\varphi^{\circ})\frac{d\KFc}{d\KRc}}_{\le0,\ \ \KFc\downarrow}
  \;-\;\underbrace{\int_0^{\KFc}\FKR(x;\KRc)\,dx}_{>0}\;<\;0 ,
\]
both terms working the same way ($\KFc\downarrow$ in $\KRc$ by Step 2, and $\FKR>0$), so
$\E[\fossil^{c}]$ strictly decreases and $\Delta^{c}$ strictly increases in $\KRc$. As
$\KRc\to\infty$, $F(x;\KRc)\to1$ for each fixed $x>0$ (renewables swamp demand), so by
dominated convergence $\E[\fossil^{c}]\to0$ and $\Delta^{c}\to\E[\fossil^{\circ}]$.
\end{proof}

\begin{remark}[The contract incentivizes renewable build-out]
Proposition~\ref{prop:saving} says the surplus contract does not merely fail to add
fossil---it actively pulls the optimal renewable capacity up, and in the cheap-renewable
regime all the way to full decarbonization. The flexible load, correctly contracted,
behaves like a subsidy for renewables financed by selling surplus power to computing.
This is the positive mechanism the paper has been after: not ``computing happens to be
green,'' but ``the right contract makes computing pay for the green build-out.''
\end{remark}

%=====================================================================
\section{Numerical validation}\label{sec:num}

All claims are checked in \texttt{verify\_week3.py} against adaptive quadrature of the
profit integrals, on non-power-law Beta scale mixtures
$\tilde x=\tilde D-\KR\tilde S$ spanning symmetric, left-skewed, right-skewed, and
bimodal demand. The suite confirms, at every adopted optimum and across the family:
the surplus contract burns zero \textsc{fcl} fossil (Prop.~\ref{prop:surplusopt}); the
surplus-channel term in \eqref{eq:focR} is strictly positive at $\KRo$, giving
$\KRc>\KRo$ (Step~1); $\KFc\le\KFo$ (Step~2); the strict carbon fall
$\E[\fossil^{c}]<\E[\fossil^{\circ}]$ holds with the \emph{same sign for every shape in
the family} (Theorem~\ref{thm:decarb}); the Pareto improvement $\Pi^{c}>\Pi^{\circ}$
with blackout held equal, blackout$^{c}=$blackout$^{\circ}=1-\varphi^{\circ}$
(Cor.~\ref{cor:triplewin}); and the exact saving \eqref{eq:exactsave} matches the direct
quadrature. As a contrast the suite also
reports the Week-2 naive-rule fossil at the same primitives, exhibiting the floor that
the contract removes. The suite reports a single pass count.

\end{document}
